% !PRINTED 453
\weberpage{453}
%<<CONT
\weq{15}{p\,\omega, \quad \frac{\omega+c}{p}}
keine zwei äquivalente, so gehören zu jedem Werte von $u$ $[p-(D,p)]$ Werte $v$, und wenn wir die Grade von (2) und (3) mit $h$ und $h'$ bezeichnen, so ergibt sich (in Übereinstimmung mit den Formeln \S\,100) folgende Beziehung:

\weq{16}{h' = [p-(D,p)]\,h.}

Nach \S\,122 sind $h$ und $h'$ die Klassenzahlen der Diskriminanten $D, D'$.

Es handelt sich also noch um die Frage, ob unter den Größen (15) nach Ausschluß der wegen (13) wegfallenden, noch äquivalente vorkommen können. Ist zunächst $p\omega$ äquivalent mit $\frac{\omega+c}{p}$, so ist:

\weq{17}{\frac{\omega+c}{p} = \frac{\gamma+\delta p\omega}{\alpha+\beta p\omega}, \quad \alpha\delta-\beta\gamma=1,}

oder:

\weq{18}{\beta p\,\omega^2 + (\alpha+\beta c p - \delta p^2)\omega + c\alpha - \gamma p = 0.}

Diese Gleichung muß aus (9) durch Multiplikation mit einer ganzen Zahl $g$ hervorgehen, und es folgt also:
\[
\beta p = gA, \quad \alpha+\beta c p-\delta p^2 = gB, \quad c\alpha - gp = gC. % [sic] printed g p; elimination from (18) requires \gamma p
\]
Da $A$ durch $p$ unteilbar angenommen war, so folgt hieraus, daß $g$ und folglich $\alpha$ durch $p$ teilbar ist. Setzen wir $g=px$, $\alpha=p\alpha'$, so folgt:

\weq{19}{\beta=Ax, \quad \alpha'+\beta c-\delta p=Bx, \quad c\alpha'-\gamma=Cx.}

Setzen wir noch

\weq{20}{\alpha'-\beta c+\delta p=y,}

so ergibt sich:

\weq{21}{\begin{aligned} 2\alpha' &= Bx+y,\\ 2(\beta c-\delta p) &= Bx-y, \end{aligned}}

und aus

\weq{22}{p\alpha'\delta-\beta\gamma=1}

folgt sodann

\weq{23}{y^2-Dx^2=4.}

Der Wert $x=0$ ist auszuschließen, weil sonst $\beta=0$ sein müßte, was der Gleichung (22) widerspricht, und wir können unbeschadet der Allgemeinheit $x$ positiv annehmen. Da überdies $D$ nach dem Modul 4 entweder $\equiv 0$ oder $\equiv 1$ ist, so
%CONT>>
